\chapter{方案设计：访问控制树到 LSSS 的 CP-ABE}

\section{总体设计思路}
本系统围绕“策略表达—矩阵转换—密钥生成—加密—解密验证”这一主线设计。加密端接收访问控制树而不是直接接收矩阵，原因在于访问树更接近用户对权限条件的自然表达；系统内部再将访问树转换为 LSSS 矩阵 $(M,\rho)$，使后续加密和解密都可以在统一的线性代数框架中进行。这样设计后，策略语义不再分散在多个手写判断分支中，而是集中体现在矩阵行、属性映射和重构方程之间的关系上。\par
从算法流程看，Setup 负责生成系统级公共参数和主密钥；KeyGen 根据用户属性集合生成私钥分量；Encrypt 根据访问树构造 LSSS 矩阵，并为矩阵每一行生成秘密份额、盲化份额和密文分量；Decrypt 则根据用户私钥中的属性集合筛选可用矩阵行，求解 $M_I^T\omega=e_1$。如果该线性方程有解，说明用户属性集合能够线性恢复秘密，系统继续恢复消息指数并解开对称载荷；如果无解，系统直接拒绝解密。因此，本方案的正确性判据可以概括为：属性集合满足访问结构，当且仅当其对应行集合能够重构向量 $(1,0,\dots,0)$。\par
系统结构如图 \ref{fig:system-architecture} 所示。CLI 层只负责参数解析、路径处理和文件读写，核心模块集中实现模运算、随机数、对称加密、策略校验、LSSS 构造、重构求解和 CP-ABE 四个主算法。这样的分层使命令行接口保持简单，也使报告中的公式、代码片段和实验文件可以逐一对应。命令行阶段输出与 JSON 文件中的策略树、矩阵 $M$、映射 $\rho$ 和密文分量共同构成了可核对的实验链路。\par
在数据组织上，各阶段之间统一通过 JSON 文件传递结果。公共参数、主密钥、用户私钥和密文分别保存，避免所有状态隐藏在一次性运行的内存变量中。密文中同时保留原始策略、转换后的 LSSS 矩阵和对称加密载荷，使实验结果具有可追踪性：当某个属性集合解密失败时，可以回到矩阵行集合和高斯消元过程解释失败原因；当解密成功时，也可以说明成功并非来自简单字符串匹配，而是来自重构系数对矩阵行的线性组合。\par
这一总体流程参考了 Bethencourt 等人~[2]和 Waters~[3]对 CP-ABE 的四算法划分。基础实现采用指数抽象模型保存群元素指数，便于观察 LSSS-CP-ABE 的核心数据流和代数抵消关系；真实 pairing 路线则在同一 \verb|Project| 中通过 \verb|-Adapter pairing| 调用 BLS12-381 后端，把公共参数、私钥和密文分量落到真实 $G_1/G_2/G_T$ 群元素上。两条路线共享访问树表示、LSSS 矩阵构造和重构系数求解思想，差异主要集中在底层群元素、配对、哈希到群、序列化和会话密钥派生方式上。
\begin{figure}[H]
\centering
\includegraphics[width=0.95\textwidth]{system_architecture.png}
\caption{系统模块结构图}
\label{fig:system-architecture}
\end{figure}

\section{访问控制树到 LSSS 的构造}

设访问控制树的每个内部节点为阈值门 $k$-of-$d$，对子节点编号为 $1,2,\dots,d$。在每个阈值门处采用 Shamir 风格的多项式分享：为该节点选择次数 $k-1$ 的多项式 $q(x)$，使得该节点持有的份额为 $q(0)$，并将子节点份额设为 $q(1),\dots,q(d)$。

由于多项式求值在域 $\mathbb{Z}_p$ 上是线性的，最终每个叶子份额都可以表示为根秘密 $s$ 与一组随机变量的线性组合。因此可以把每个叶子份额写成：

\[
    \lambda_i = M_i \cdot \mathbf{v},
\]

其中 $\mathbf{v}=(s,r_2,\dots,r_n)$，$M_i$ 为第 $i$ 个叶子对应的系数向量。收集所有叶子的系数向量即可得到 LSSS 矩阵 $M$，并用 $\rho(i)$ 记录第 $i$ 行对应的属性标签。

本项目采用如下列数计算方式：

\[
    n = 1+\sum_{\text{gate }v}(k_v-1).
\]

其中第一列对应根秘密 $s$，每个阈值门额外贡献 $k_v-1$ 个随机维度。对于示例策略 $(A\wedge B)\vee C$，根 OR 门为 $1$-of-$2$，不引入随机维度；内部 AND 门为 $2$-of-$2$，引入一个随机维度。因此最终矩阵列数为 $2$。

这一列数设计体现了“按需引入随机性”的原则。OR 门只需要把父节点份额复制到某个满足的分支，不需要额外随机变量；AND 和一般阈值门则需要通过多项式分享隐藏父节点份额，因此会增加随机维度。这样得到的矩阵既能表达策略，又不会无谓增加列数。

该设计与 Liu、Cao 和 Wong~[4]“直接处理阈值访问树”的思想一致：每个内部节点均作为阈值访问结构参与矩阵生成，而不需要先把阈值策略展开成大量 AND/OR 组合。对于课程题目要求的“使用访问控制树构建 LSSS 矩阵”，这一处理方式使策略树的分支数、层次结构和阈值参数能够直接进入矩阵构造过程，从而区别于只面向固定样例或二叉 AND/OR 公式的实现。

\subsection{访问策略树一般性的体现}

本文所称访问策略树的一般性，主要限定在单调访问结构范围内。也就是说，策略树由属性叶子和阈值门组成，不引入 NOT 这类非单调条件；在这一范围内，系统并不依赖某个固定样例或固定层数的树形结构。策略节点只需要满足两类基本形式：叶子节点保存属性名，内部节点保存阈值 $k$ 与子节点列表。只要 $1\leq k\leq d$，其中 $d$ 为子节点个数，该内部节点就可以作为根节点，也可以出现在任意子树位置；其子节点既可以是属性叶子，也可以继续是阈值门。因此，AND、OR、$2$-of-$3$ 以及更深层的嵌套阈值结构都可以落在同一种访问树表示之下。

这种一般性进一步体现在 LSSS 矩阵构造过程中。转换算法对每个内部节点采用相同的阈值门处理规则：根据节点阈值 $k$ 引入 $k-1$ 个随机维度，以父节点份额作为多项式常数项，再把对子节点编号 $1,\dots,d$ 的求值结果继续向下传递。叶子节点只负责把最终份额向量写入矩阵行，并通过 $\rho$ 记录对应属性。由于递归过程不需要区分“这是 AND 门”“这是 OR 门”或“这是第几层节点”，树的分支数、嵌套深度和阈值取值都由输入策略本身决定。

解密阶段同样没有回到原始访问树上逐个门进行布尔判断，而是根据用户属性集合筛选满足 $\rho(i)\in S$ 的矩阵行，并求解 $M_I^T\omega=e_1$。如果存在重构系数 $\omega$，说明这些行能够恢复根秘密对应的第一维；如果不存在，则属性集合不满足策略。由此，访问策略的满足性被统一转化为线性重构问题。第五章中 $2$-of-$3(D,E,F)$ 的成功与失败测试、重复属性 $X\wedge X$ 测试以及非法策略输入测试，分别从一般阈值门、多行映射同一属性和输入边界三个角度验证了这种设计不是针对示例策略 $(A\wedge B)\vee C$ 的硬编码。

\section{CP-ABE 四个核心算法}

系统包含四个基本算法：Setup、KeyGen、Encrypt 和 Decrypt。它们在命令行中分别对应 \verb|setup|、\verb|keygen|、\verb|encrypt| 和 \verb|decrypt| 四个子命令。

\subsection{Setup}

Setup 生成系统公共参数和主密钥。真实 CP-ABE 中会选择双线性群、生成元和指数。本实验在指数抽象模型下固定素数域 $\mathbb{Z}_p$，并随机选择主密钥指数 $\alpha$ 和 $a$：

\[
    PK=(p,g,g^a,e(g,g)^\alpha),\quad MSK=(\alpha,a).
\]

在 JSON 文件中，$g$ 以指数 $1$ 表示，$g^a$ 存储为 $a$，$e(g,g)^\alpha$ 存储为 $\alpha$。

在真实 pairing 路线中，Setup 使用 BLS12-381 的曲线阶作为标量域，公共参数保存 $g_1\in G_1$、$g_2\in G_2$、$g_1^a$ 以及 $e(g_2,g_1)^\alpha\in G_T$：

\[
    PK=(p,g_1,g_2,g_1^a,e(g_2,g_1)^\alpha),\quad MSK=(\alpha,a).
\]

这些群元素写入 JSON 时会被编码为坐标或系数数组，而不是单个指数值。

Setup 阶段的设计目标是提供后续三个算法共同依赖的公共基础。公共参数中的 $a$ 和 $\alpha$ 会分别进入密文分量和私钥分量，解密时再通过线性组合抵消。主密钥只在 KeyGen 阶段使用，不写入密文，也不参与用户解密命令，这与真实 ABE 系统中“授权中心持有主密钥”的角色划分一致。

\subsection{KeyGen}

KeyGen 为用户属性集合 $S$ 生成私钥。私钥由用户级分量和属性级分量两部分组成：$K$ 与 $L$ 是该用户私钥的全局分量，$K_x$ 是属性 $x\in S$ 对应的分量。这里的下标 $x$ 表示属性索引，因此一个属性集合会对应多个 $K_x$；为避免符号过于繁琐，公式中省略了用户下标，但不同用户重新运行 KeyGen 时会得到不同的具体数值。算法随机选取用户独立随机数 $t$，并生成：

\[
    K=g^\alpha g^{at},\quad L=g^t,\quad K_x=H(x)^t.
\]

上式是论文中常见的群元素写法。本项目采用指数抽象模型保存同一组私钥分量，因此不是再生成一套新的 $K,L,K_x$，而是把这些群元素对应的指数记录下来：

\[
    K=\alpha+at,\quad L=t,\quad K_x=H(x)t.
\]

同一个用户所有属性分量共享同一个 $t$，这体现了抗合谋设计思想：不同用户即使合并属性，也难以把不同随机数 $t$ 的私钥分量拼成一个合法用户的私钥结构。

在本实验实现中，抗合谋性质主要通过公式结构体现，而不是通过完整安全证明体现。报告中保留这一点，是为了说明 KeyGen 不只是为每个属性单独生成一个哈希值，而是把用户级随机数 $t$ 同时绑定到 $K$、$L$ 和所有 $K_x$ 分量中。后续解密公式中的抵消正是依赖这些分量共享同一随机数。

真实 pairing 版沿用这一结构，但分量所在群不同：$K\in G_1$，$L\in G_2$，$K_x\in G_1$。属性哈希 $H(x)$ 由 \verb|hash_to_G1| 产生，使属性名直接映射为 $G_1$ 群元素。该处理比“哈希为标量再写入指数模型”更接近真实配对系统中的属性映射方式。

\subsection{Encrypt}

Encrypt 输入公共参数、访问策略和明文。首先构造 LSSS 矩阵 $(M,\rho)$，然后选择秘密 $s$ 和两个向量：

\[
    \mathbf{v}=(s,r_2,\dots,r_n),\quad
    \mathbf{u}=(0,z_2,\dots,z_n).
\]

对每一行计算：

\[
    \lambda_i=M_i\cdot \mathbf{v},\quad
    w_i=M_i\cdot \mathbf{u}.
\]

其中 $\lambda_i$ 是秘密份额，$w_i$ 是盲化份额。$\mathbf{u}$ 的第一维为 $0$，使得满足策略的线性组合会把盲化项重构为 $0$，从而在解密中抵消。

密文核心分量为：

\[
    C=m+\alpha s,\quad C_0=s,
\]

以及每一行对应的：

\[
    C_i=a\lambda_i-H(\rho(i))w_i,\quad D_i=w_i.
\]

这里的加法、乘法和减法均在模 $p$ 下进行。

真实 pairing 版中，受 CP-ABE 保护的对象是随机 $G_T$ 会话元素 $M$。密文分量对应为：

\[
    C=M\cdot e(g_2,g_1)^{\alpha s},\quad C_0=g_2^s,
\]

以及每一行对应的：

\[
    C_i=(g_1^a)^{\lambda_i}\cdot H(\rho(i))^{-w_i},\quad D_i=g_2^{w_i}.
\]

因此，指数版中的 $C,C_0,C_i,D_i$ 可以看作真实群元素公式在指数层面的投影；pairing 版则直接保存 $G_T$、$G_2$、$G_1$、$G_2$ 中的实际元素。

Encrypt 阶段的两个向量 $\mathbf{v}$ 和 $\mathbf{u}$ 分别承担不同作用。$\mathbf{v}$ 负责分享真正的秘密 $s$，决定满足策略时能恢复出 $\alpha s$；$\mathbf{u}$ 的首项为 $0$，负责生成可抵消的盲化项，防止单个密文行直接泄露份额结构。只有当用户属性集合能够组合出 $(1,0,\dots,0)$ 时，盲化项才会在线性重构中同时消去。

实验三课件~[1]中 Waters 2011 加密步骤强调对访问矩阵每一行计算 $\lambda_i=\mathbf{v}\cdot M_i$；Waters~[3]则把密文与 LSSS 访问结构 $(M,\rho)$ 绑定。本项目在此基础上额外保留 $\mathbf{u}$ 盲化向量，是为了在指数抽象模型中呈现解密抵消过程：满足策略时 $\mathbf{v}$ 的第一维恢复 $s$，而 $\mathbf{u}$ 的第一维为 $0$，因此盲化项随同线性重构被消去。

\subsection{Decrypt}

Decrypt 输入公共参数、用户私钥和密文。首先筛选所有满足 $\rho(i)\in S$ 的行，构造行集合 $I$。然后求解：

\[
    M_I^T\omega=e_1,
\]

其中 $e_1=(1,0,\dots,0)^T$。如果该方程无解，说明用户属性集合不能满足策略，算法抛出错误并拒绝解密。

如果方程有解，则计算：

\[
\begin{aligned}
A &= \sum_{i\in I}\omega_i(C_iL+D_iK_{\rho(i)})\\
  &= \sum_{i\in I}\omega_i((a\lambda_i-H(\rho(i))w_i)t+w_iH(\rho(i))t)\\
  &= at\sum_{i\in I}\omega_i\lambda_i\\
  &= ats.
\end{aligned}
\]

再计算：

\[
    B=C_0K-A=s(\alpha+at)-ats=\alpha s,
\]

最终恢复：

\[
    m=C-B.
\]

恢复出的 $m$ 用于派生对称密钥，进而验证 HMAC 并解密 AES-CBC 载荷。

在真实 pairing 版中，解密端计算：

\[
\prod_{i\in I}\left(e(L,C_i)\cdot e(D_i,K_{\rho(i)})\right)^{\omega_i}
    = e(g_2,g_1)^{ats},
\]

再用

\[
    e(C_0,K)=e(g_2,g_1)^{\alpha s+ats}
\]

消去 $ats$，得到 $e(g_2,g_1)^{\alpha s}$。随后从 $C$ 中除去该遮蔽项即可恢复 $G_T$ 会话元素，并由该元素派生 AES/HMAC 密钥。

Decrypt 阶段的失败路径同样重要。如果求不出 $\omega$，程序不会尝试“部分解密”，而是立即抛出不满足策略的错误。这使访问控制语义和线性方程求解结果保持一致：有解代表满足策略，无解代表拒绝访问。后续实验中的 $\{A\}$ 和 $\{D\}$ 两个失败场景都用于验证这一点。

Waters~[3]对 Decrypt 的描述中，满足策略的属性集合对应行集合 $I=\{i:\rho(i)\in S\}$，并存在一组常数使得这些行可以线性恢复秘密。实现时选择直接求解 $M_I^T\omega=e_1$，能够把这一理论条件变成可执行的判定过程。这也是本项目没有单独写树满足性递归判断的原因：矩阵可解性已经完整表达了访问结构语义。

\section{数据流设计}

图 \ref{fig:encrypt-decrypt-flow} 展示了系统的主要数据流。加密端从策略 JSON 和明文出发，生成密文 JSON；解密端从密文 JSON 和用户私钥出发，先求解 LSSS 重构系数，再恢复对称密钥并输出明文。

数据流设计中特别保留了中间结构：密文不仅保存对称密文，还保存策略树、矩阵 $M$ 和映射 $\rho$。这样做会使密文 JSON 更大，但能够把“访问树已经被转换为矩阵”这一过程固化在实验文件中，也使失败场景能够回到矩阵行和属性映射层面分析原因。

\begin{figure}[H]
\centering
\includegraphics[width=0.95\textwidth]{encrypt_decrypt_flow.png}
\caption{加密与解密数据流}
\label{fig:encrypt-decrypt-flow}
\end{figure}

\section{工程化接口设计}

为便于实验展示与复现，系统提供命令行接口：

\begin{table}[H]
\centering
\caption{命令行接口}
\begin{tabular}{p{0.18\textwidth} p{0.70\textwidth}}
\toprule
命令 & 作用 \\
\midrule
\verb|setup| & 生成公共参数 \verb|public.json| 和主密钥 \verb|master.json|。\\
\verb|keygen| & 输入公共参数、主密钥和属性集合，生成用户私钥 JSON。\\
\verb|encrypt| & 输入公共参数、策略 JSON 和明文文件，输出密文 JSON。\\
\verb|decrypt| & 输入公共参数、用户私钥和密文 JSON，输出恢复后的明文文件。\\
\verb|-Adapter| \verb|pairing| & 启用 BLS12-381 真实 pairing 后端。\\
\bottomrule
\end{tabular}
\end{table}

这种 CLI 设计使每个算法阶段都保留单独的文件结果，便于在实验分析中追踪 Setup、KeyGen、Encrypt 和 Decrypt 的状态变化。JSON 格式虽然不如二进制编码紧凑，但具有可读性强、便于调试、便于展示矩阵和中间分量的优点。

与只提供函数接口相比，CLI 使实验过程具有更清晰的阶段边界。逐条命令和烟雾测试脚本都会在 \path{tests/artifacts} 目录下留下公共参数、主密钥、用户私钥、密文和解密输出，这些文件与报告中的字段说明可以相互对应。

\section{安全性边界与实现取舍}

本节的安全性边界主要用于明确本实验设计所覆盖的问题范围。本文实现的核心目标是验证“访问控制树能够被转换为 LSSS 矩阵，并进一步嵌入 CP-ABE 四算法流程”这一主线。因此，报告中的安全性讨论并不等同于完整 CP-ABE 论文中的安全证明，而是围绕访问结构正确性、代数重构正确性、密文载荷完整性和实现边界展开。这样处理可以避免把“算法流程能够正确运行”和“方案已经满足标准模型或选择明文安全证明”混为一谈。

在指数抽象模型中，真实群元素被表示为指数，群乘法、幂运算和双线性对分别对应指数加法、指数乘法和指数乘积。该模型保留了 CP-ABE 解密公式中最关键的代数关系：满足访问结构的属性集合能够通过重构系数恢复秘密 $s$，并使盲化项在线性组合中抵消；不满足访问结构的属性集合无法求得 $M_I^T\omega=e_1$ 的解，因此不能进入消息恢复流程。真实 pairing 路线进一步把这些关系放入 BLS12-381 的群元素和配对运算中执行，能够检验群元素序列化、配对调用和 $G_T$ 密钥派生是否与访问结构逻辑正确衔接。

同时，即使接入真实 pairing 库，报告中的安全性讨论仍不能等同于完整论文安全证明。真实 CP-ABE 方案中的安全性通常依赖素数阶配对友好曲线、随机预言机或标准模型假设、群元素序列化约束、严格随机数生成过程以及与论文完全一致的算法分量。真实 pairing 路线已经覆盖真实群元素、哈希到群、配对和基本反序列化检查，但仍需把“可运行的真实群实现”和“满足某篇论文安全模型的完整实现”区分开。

\begin{table}[H]
\centering
\caption{本实验实现的安全性边界}
\begin{tabular}{>{\raggedright\arraybackslash}p{0.20\textwidth} >{\raggedright\arraybackslash}p{0.35\textwidth} >{\raggedright\arraybackslash}p{0.33\textwidth}}
\toprule
分析对象 & 本实验覆盖内容 & 边界说明 \\
\midrule
访问策略 & 支持属性叶子、AND/OR 和一般阈值门，并由访问树生成 LSSS 矩阵 & 限定于单调访问结构，不包含 NOT 等非单调条件 \\
解密判定 & 通过 $M_I^T\omega=e_1$ 判断属性集合是否满足策略 & 正确性依赖矩阵构造与有限域线性求解实现 \\
代数抵消 & 在指数抽象模型中验证 $\lambda_i$、$w_i$、$K$、$L$、$K_x$ 的抵消关系，并在 pairing 版中用真实配对复核同一访问逻辑 & pairing 版提升了实现真实性，但仍不等同于形式化安全证明 \\
载荷保护 & 使用 AES-CBC 加密明文，并用 HMAC-SHA256 验证 IV 与密文完整性 & 对称层保护文件载荷，不改变 CP-ABE 访问结构语义 \\
输入边界 & 对非法策略、重复属性和特殊属性名进行测试 & 工程健壮性测试不等同于形式化安全证明 \\
\bottomrule
\end{tabular}
\end{table}

从实现取舍看，本文先建立了相对独立的群运算抽象层，再在同一 \verb|Project| 中接入真实 BLS12-381 后端。这样做使访问树解析、LSSS 矩阵生成、有限域高斯消元、密文 JSON 结构和 CLI 流程能够先被完整验证，再逐步替换底层群元素生成、群运算、配对、哈希到群和序列化部分。实际接入结果也说明，访问树到矩阵、行集合筛选和重构系数求解等逻辑可以保持稳定，变化最大的部分集中在群元素表示和密文/私钥 JSON 编码上。

这一边界划分也决定了实验结果分析的侧重点。对于本项目，关键分析对象不是群元素计算性能或安全证明强度，而是策略语义是否被正确转换为矩阵，满足策略的属性集合是否能够恢复消息指数，不满足策略的属性集合是否被拒绝，以及密文载荷被篡改时是否无法输出明文。第五章设计的成功解密、失败解密、一般阈值门、重复属性、非法策略和密文篡改测试，正是围绕这些边界进行验证。

Riepel 和 Wee 提出的 FABEO 方案~[5]表明，真实 ABE 方案还会进一步关注自适应安全、多密文安全、配对次数、群元素大小以及属性多次出现带来的效率与安全权衡。与这些研究目标相比，本文的方案设计更侧重课程实验要求中的访问控制树、LSSS 矩阵和 CP-ABE 基本流程。真实 pairing 路线提高了底层密码运算的真实性，但本文仍把结论限定在可由实现和实验直接支撑的范围内。
